\textbf{Phần 1. Trắc nghiệm lựa chọn}
\setcounterpageref{ex}{0}

\begin{ex} %[5H3N18]
Số đo thích hợp nhất chỉ diện tích của một trường tiểu học là
\choice
{$1{,}5\text{ m}^2$}
{$1{,}5\text{ dm}^2$}
{\True $1{,}5\text{ ha}$}
{$1{,}5\text{ cm}^2$}
\loigiai{
Diện tích khuôn viên một trường học thường được đo bằng héc-ta, phù hợp nhất là $1{,}5\text{ ha}$ ($15\,000\text{ m}^2$).\\
\textbf{Đáp án của câu hỏi là C}
}
\end{ex}

\begin{ex} %[5H3N18]
Số đo thích hợp chỉ diện tích của một bức tranh treo tường là
\choice
{$1\text{ km}^2$}
{$1\text{ ha}$}
{$1\text{ mm}^2$}
{\True $1\text{ m}^2$}
\loigiai{
Diện tích của một bức tranh treo tường trong phòng phù hợp nhất là $1\text{ m}^2$.\\
\textbf{Đáp án của câu hỏi là D}
}
\end{ex}

\begin{ex} %[5H3H18]
Đổi $25\,000\text{ m}^2$ ra héc-ta ta được kết quả là
\choice
{$0{,}25\text{ ha}$}
{\True $2{,}5\text{ ha}$}
{$25\text{ ha}$}
{$250\text{ ha}$}
\loigiai{
Ta có: $25\,000\text{ m}^2 = 25\,000 : 10\,000 = 2{,}5\text{ ha}$.\\
\textbf{Đáp án của câu hỏi là B}
}
\end{ex}

\begin{ex} %[5H3H18]
Đổi $7{,}5\text{ km}^2$ ra héc-ta ta được
\choice
{\True $750\text{ ha}$}
{$75\text{ ha}$}
{$7\,500\text{ ha}$}
{$75\,000\text{ ha}$}
\loigiai{
Ta có: $7{,}5\text{ km}^2 = 7{,}5 \times 100 = 750\text{ ha}$.\\
\textbf{Đáp án của câu hỏi là A}
}
\end{ex}

\begin{ex} %[5H3H18]
Đổi $0{,}48\text{ km}^2$ ra héc-ta ta được
\choice
{$4{,}8\text{ ha}$}
{\True $48\text{ ha}$}
{$480\text{ ha}$}
{$4\,800\text{ ha}$}
\loigiai{
Ta có: $0{,}48\text{ km}^2 = 0{,}48 \times 100 = 48\text{ ha}$.\\
\textbf{Đáp án của câu hỏi là B}
}
\end{ex}

\begin{ex} %[5H3H18]
Điền dấu thích hợp vào chỗ chấm: $4\text{ m}^2\ 2\text{ dm}^2 \dots 42\text{ m}^2$
\choice
{$>$}
{$=$}
{\True $<$}
{Không so sánh được}
\loigiai{
Đổi: $4\text{ m}^2\ 2\text{ dm}^2 = 4{,}02\text{ m}^2 < 42\text{ m}^2$.\\
\textbf{Đáp án của câu hỏi là C}
}
\end{ex}

\begin{ex} %[5H3H18]
Điền dấu thích hợp vào chỗ chấm: $2\,800\text{ ha} \dots 28\text{ km}^2$
\choice
{$>$}
{\True $=$}
{$<$}
{Không so sánh được}
\loigiai{
Đổi $28\text{ km}^2 = 28 \times 100 = 2\,800\text{ ha}$. Vậy $2\,800\text{ ha} = 28\text{ km}^2$.\\
\textbf{Đáp án của câu hỏi là B}
}
\end{ex}

\begin{ex} %[5H3H18]
Đổi $450\text{ dm}^2$ ra mét vuông ta được số thập phân là
\choice
{\True $4{,}5\text{ m}^2$}
{$0{,}45\text{ m}^2$}
{$45\text{ m}^2$}
{$450\text{ m}^2$}
\loigiai{
Ta có: $450\text{ dm}^2 = 450 : 100 = 4{,}5\text{ m}^2$.\\
\textbf{Đáp án của câu hỏi là A}
}
\end{ex}

\begin{ex} %[5H3V18]
Người ta muốn xây dựng khu vui chơi kết hợp khu cắm trại trên mảnh đất $2\text{ ha}$. Biết khu vui chơi chiếm $\dfrac{3}{4}$ diện tích mảnh đất. Diện tích khu cắm trại là
\choice
{$15\,000\text{ m}^2$}
{$10\,000\text{ m}^2$}
{\True $5\,000\text{ m}^2$}
{$7\,500\text{ m}^2$}
\loigiai{
Đổi $2\text{ ha} = 20\,000\text{ m}^2$.\\
Khu cắm trại chiếm: $1 - \dfrac{3}{4} = \dfrac{1}{4}$ diện tích mảnh đất.\\
Diện tích khu cắm trại là: $20\,000 \times \dfrac{1}{4} = 5\,000\text{ (m}^2\text{)}$.\\
\textbf{Đáp án của câu hỏi là C}
}
\end{ex}

\begin{ex} %[5H3V18]
Bác Ba muốn dùng gỗ để lát sàn phòng khách hình chữ nhật có chiều dài $6\text{ m}$, chiều rộng $5\text{ m}$. Loại gỗ bác chọn có giá $300\,000$ đồng/$1\text{ m}^2$. Số tiền bác Ba cần trả là
\choice
{$3\,000\,000$ đồng}
{$6\,000\,000$ đồng}
{$12\,000\,000$ đồng}
{\True $9\,000\,000$ đồng}
\loigiai{
Diện tích sàn phòng khách là: $6 \times 5 = 30\text{ (m}^2\text{)}$.\\
Số tiền bác Ba cần trả là: $30 \times 300\,000 = 9\,000\,000$ (đồng).\\
\textbf{Đáp án của câu hỏi là D}
}
\end{ex}

\begin{ex} %[5H3V18]
Người ta cắt một tấm kính hình chữ nhật thành $5$ tấm kính nhỏ giống nhau để đóng khung ảnh, mỗi tấm nhỏ có kích thước $6\text{ dm} \times 4\text{ dm}$. Diện tích tấm kính ban đầu theo đơn vị mét vuông là
\choice
{\True $1{,}2\text{ m}^2$}
{$12\text{ m}^2$}
{$0{,}12\text{ m}^2$}
{$24\text{ m}^2$}
\loigiai{
Diện tích một tấm kính nhỏ là: $6 \times 4 = 24\text{ (dm}^2\text{)}$.\\
Diện tích tấm kính ban đầu là: $24 \times 5 = 120\text{ (dm}^2\text{)} = 1{,}2\text{ m}^2$.\\
\textbf{Đáp án của câu hỏi là A}
}
\end{ex}

\begin{ex} %[5H3V18]
Một mảnh đất hình chữ nhật gồm $12$ ô đất hình vuông, mỗi ô có cạnh $10\text{ m}$. Người ta dùng mảnh đất đó để xây nhà máy. Diện tích xây nhà máy là bao nhiêu héc-ta?
\choice
{$1{,}2\text{ ha}$}
{$12\text{ ha}$}
{$0{,}012\text{ ha}$}
{\True $0{,}12\text{ ha}$}
\loigiai{
Diện tích $1$ ô vuông là: $10 \times 10 = 100\text{ (m}^2\text{)}$.\\
Diện tích mảnh đất là: $12 \times 100 = 1\,200\text{ (m}^2\text{)}$.\\
Đổi $1\,200\text{ m}^2 = 1\,200 : 10\,000 = 0{,}12\text{ ha}$.\\
\textbf{Đáp án của câu hỏi là D}
}
\end{ex}

\textbf{Phần 2. Trắc nghiệm đúng sai}
\setcounterpageref{ex}{0}

\begin{ex} %[5H3H18]
Xét tính đúng sai của các khẳng định sau về các phép đổi đơn vị đo diện tích:
\choiceTF
{\True $412\,000\text{ m}^2 = 41{,}2\text{ ha}$.}
{$8\,000\text{ m}^2 = 8\text{ ha}$.}
{\True $11\text{ km}^2 = 1\,100\text{ ha}$.}
{\True $8\text{ ha} = 80\,000\text{ m}^2$.}
\loigiai{
\begin{itemize}
    \item $412\,000\text{ m}^2 = 412\,000 : 10\,000 = 41{,}2\text{ ha}$.
    \item $8\,000\text{ m}^2 = 8\,000 : 10\,000 = 0{,}8\text{ ha} \neq 8\text{ ha}$.
    \item $11\text{ km}^2 = 11 \times 100 = 1\,100\text{ ha}$.
    \item $8\text{ ha} = 8 \times 10\,000 = 80\,000\text{ m}^2$.
\end{itemize}
\begin{itemize}
    \item a) \textbf{Đúng}.
    \item b) \textbf{Sai}: Kết quả là $0{,}8\text{ ha}$.
    \item c) \textbf{Đúng}.
    \item d) \textbf{Đúng}.
\end{itemize}
}
\end{ex}

\begin{ex} %[5H3V18]
Một mảnh đất có diện tích $2\text{ ha}$ được chia thành hai phần: khu vui chơi chiếm $\dfrac{3}{4}$ diện tích và khu cắm trại chiếm $\dfrac{1}{4}$ diện tích. Xét tính đúng sai của các khẳng định sau:
\choiceTF
{\True Diện tích cả mảnh đất là $20\,000\text{ m}^2$.}
{Diện tích khu vui chơi là $1\,500\text{ m}^2$.}
{\True Diện tích khu cắm trại là $5\,000\text{ m}^2$.}
{\True Diện tích khu vui chơi lớn hơn diện tích khu cắm trại là $1\text{ ha}$.}
\loigiai{
\begin{itemize}
    \item $2\text{ ha} = 20\,000\text{ m}^2$.
    \item Diện tích khu vui chơi: $20\,000 \times \dfrac{3}{4} = 15\,000\text{ m}^2 \neq 1\,500\text{ m}^2$.
    \item Diện tích khu cắm trại: $20\,000 \times \dfrac{1}{4} = 5\,000\text{ m}^2$.
    \item Hiệu diện tích: $15\,000 - 5\,000 = 10\,000\text{ m}^2 = 1\text{ ha}$.
\end{itemize}
\begin{itemize}
    \item a) \textbf{Đúng}.
    \item b) \textbf{Sai}: Diện tích là $15\,000\text{ m}^2$.
    \item c) \textbf{Đúng}.
    \item d) \textbf{Đúng}.
\end{itemize}
}
\end{ex}

\textbf{Phần 3. Trả lời ngắn}
\setcounterpageref{ex}{0}

\begin{ex} %[5H3H18]
Đổi $8\,000\text{ m}^2$ ra héc-ta. (Nhập kết quả dưới dạng số thập phân).
\shortans{0,8}
\loigiai{
\textbf{Đáp số:} $0{,}8$\\
$8\,000\text{ m}^2 = 8\,000 : 10\,000 = 0{,}8\text{ ha}$.
}
\end{ex}

\begin{ex} %[5H3H18]
Đổi $0{,}48\text{ km}^2$ ra héc-ta.
\shortans{48}
\loigiai{
\textbf{Đáp số:} $48$\\
$0{,}48\text{ km}^2 = 0{,}48 \times 100 = 48\text{ ha}$.
}
\end{ex}

\begin{ex} %[5H3H18]
Đổi $450\text{ dm}^2$ ra mét vuông. (Nhập kết quả dưới dạng số thập phân).
\shortans{4,5}
\loigiai{
\textbf{Đáp số:} $4{,}5$\\
$450\text{ dm}^2 = 450 : 100 = 4{,}5\text{ m}^2$.
}
\end{ex}

\begin{ex} %[5H3V18]
Bác Ba lát sàn phòng khách dài $6\text{ m}$, rộng $5\text{ m}$ với giá gỗ $300\,000$ đồng/$1\text{ m}^2$. Hỏi bác Ba phải trả bao nhiêu triệu đồng tiền gỗ?
\shortans{9}
\loigiai{
\textbf{Đáp số:} $9$\\
Diện tích: $6 \times 5 = 30\text{ m}^2$. Chi phí: $30 \times 300\,000 = 9\,000\,000$ đồng $= 9$ triệu đồng.
}
\end{ex}

\begin{ex} %[5H3V18]
Cắt một tấm kính lớn thành $5$ tấm nhỏ cùng kích thước $6\text{ dm} \times 4\text{ dm}$. Tính diện tích tấm kính lớn ban đầu theo đơn vị mét vuông. (Nhập kết quả dưới dạng số thập phân).
\shortans{1,2}
\loigiai{
\textbf{Đáp số:} $1{,}2$\\
Diện tích là: $5 \times (6 \times 4) = 120\text{ dm}^2 = 1{,}2\text{ m}^2$.
}
\end{ex}

\begin{ex} %[5H3V18]
Mảnh đất gồm $12$ ô vuông nhỏ, mỗi ô có cạnh dài $10\text{ m}$. Tính diện tích mảnh đất đó theo đơn vị héc-ta. (Nhập kết quả dưới dạng số thập phân).
\shortans{0,12}
\loigiai{
\textbf{Đáp số:} $0{,}12$\\
Diện tích: $12 \times (10 \times 10) = 1\,200\text{ m}^2 = 0{,}12\text{ ha}$.
}
\end{ex}
